\scanpage{133}
\mspage{82}

que
\[
 k\in\mathcal W\Longleftrightarrow k_0\in\mathcal W
\]
d'où
\[
 i\in\mathcal W\Longleftrightarrow(k_0:\dot T\to T)\in\mathcal W.
\]
De \unsure{=|même} \deleted{\unsure{pour}} on trouve\deleted{\unsure{r}}, en appliquant ceci à
$(X',U')$ au lieu de $(X,U)$
\[
 i'\in\mathcal W\Longleftrightarrow(k'_0:\dot T'\unsure{-|\to}T')\in\mathcal W.
\]
Or on a $T\xrightarrow{\sim}T'$ (\unsure{iso|isom}), et par la prop.~4
cet iso. transforme $\dot T$ en $\dot T'$. \deleted{D'où la conclusion}\inserted{Donc}
\[
 k_0\in\mathcal W\Longleftrightarrow k'_0\in\mathcal W\quad\text{(trivialement),}
\]
d'où a).

Prouvons b). Le carré comm. induit par $(*)$ p.~80
\diagram{\begin{tikzcd}[ampersand replacement=\&]
 U \arrow[d,"p"'] \& \overline U \arrow[l,"i_0"'] \arrow[d,"{\deleted{\unsure{q_{\overline U}|p_{\overline U}}}\overline p}"]\\
 U' \arrow[r,hook,"i'_0"] \& \overline U'
\end{tikzcd}}\noindent,
où $i_0,i'_0\in\mathcal W$ par hypothèse, montre que
\[
 p\in\mathcal W\Longleftrightarrow\overline p\deleted{\unsure{q_{\overline U}|p_{\overline U}}}\in\mathcal W,
\]
donc tout revient à prouver que
\[
 \overline p\in\mathcal W\overset{?}{\Longleftrightarrow}q\in\mathcal W.
\]
Mais $q$ induit un hom du carré $(*)$ déduit
de \deleted{$(X\unsure{,|q}\overline U,$} $(X,\overline U,T)$ (\unsure{avec|car} $X=\overline U\cup T$), dans le
carré similaire pour $X'$
\diagram{\begin{tikzcd}[ampersand replacement=\&]
 \overline U'\cap T'=\dot T' \arrow[r,"k'_0"] \arrow[d,"j'_0"'] \& T' \arrow[d,"j'"]\\
 \overline U' \arrow[r,"k'"] \& X'
\end{tikzcd}}\noindent,
les flèches verticales étant notées $q_{\dot T},q_T,\overline p\deleted{\unsure{q|U}},q$.

% archive sheet number: 133
